\section{极小Cauchy滤子}

记号约定:$X$是一致空间,$\mathfrak{S}$是$X$上所有对称邻里构成的集合.

\begin{proposition}{极小Cauchy滤子的存在性和构造}{}
	\begin{enumerate}
		\item 对$X$上的每个Cauchy滤子$\mathfrak{F}$,比$\mathfrak{F}$粗的极小Cauchy滤子$\mathfrak{F}_{0}$存在且唯一.
		\item 若$\mathfrak{F}$是$X$上的Cauchy滤子,$\mathfrak{B}$是$\mathfrak{F}$的基,则$\left\{V\left(M\right)\middle|M\in\mathfrak{B},V\in\mathfrak{S}\right\}$是$\mathfrak{F}_{0}$的基.
	\end{enumerate}
\end{proposition}

\begin{corollary}{点的邻域滤子是极小Cauchy滤子}{}
	对$X$中每一点$x$,它在$X$中的邻域滤子是极小Cauchy滤子.
\end{corollary}

\begin{corollary}{柯西滤子的聚点即极限点}{}
	设$\mathfrak{F}$是$X$上的Cauchy滤子,则$\mathfrak{F}$的每个聚点都是自身的极限点.
\end{corollary}

\begin{corollary}{较粗柯西滤子的收敛性}{}
	设$x\in X$,则$X$上每个比收敛到$x$的滤子更粗的Cauchy滤子都收敛于$x$.
\end{corollary}

\begin{corollary}{极小柯西滤子的开集基}{}
	若$\mathcal{F}$是$X$上的极小Cauchy滤子,则$\mathfrak{F}$有一个基由$X$中的开集构成.
\end{corollary}